\documentclass{article}
\usepackage{amsmath,amssymb,hyperref}
\title{Derived fixed-interface state synchronization in the module completion}
\date{15 November 2026}
\begin{document}\maketitle
\paragraph{Status.} The one-input normal-lax preservation calculation and the fixed-interface application have passed independent written-proof review under the stated hypotheses. It discharges structural hypotheses using the actual derived parallel and module construction, rather than assuming arbitrary new residual arrows already support parallel.

\paragraph{One-input preservation theorem.} Let $U:Q\to R$ be normal lax and locally cocontinuous between locally cocomplete bicategories with separately cocontinuous composition. Its enriched module extension $U_*$ is normal lax. For every small normal lax diagram $F:D\to\operatorname{Mod}(Q)$, there is a canonical enriched equivalence
\[
 U_*(A_F)\simeq A_{U_*F}
\]
compatible with the native incoming lax cone. On homs it is $U(\coprod_aF(a))\cong\coprod_a UF(a)$. Multiplication uses the same directed compositor of $U$ and the same diagram compositor on either side; normality identifies units. This proof does not invert either compositor.

The projection $p_j:A_F\to Fj$ has components $\coprod_{a:i\to j}F(a)(y,x)$. These identify with the target projection after applying $U$. More strongly, $p_j$ is the restriction module along the enriched inclusion $Fj\to A_F$. For a column $W:X\to A_F$, module Yoneda gives $p_j\odot W\cong W_j$. The canonical comparison $U_*p_j\odot U_*W\to U_*(p_j\odot W)$ identifies with the identity on $U(W_j)$. Thus the actual universal cone and its modification equations match. This uses restriction exactness, not invertibility of an arbitrary module compositor. The ambient restriction principle is consistent with Shulman, \emph{Framed bicategories and monoidal fibrations}, Proposition 6.8, \url{https://arxiv.org/abs/0706.1286}; the component argument here does not assume the original pure interface was framed.

\paragraph{Fixed-interface consequence of binary assembly.} Let $T_*$ be the paired module parallel induced by the separately locally cocontinuous original normal-lax parallel. Fix an enriched category $B$. Taking the second diagram to be the constant normal diagram on $B$ indexed by the terminal category, the product-index theorem gives
\[
 T_*(A_F,B)\simeq A_{T_*(F,B)}.
\]
Here $B$ can have many objects; its full identity module, not merely an identity arrow at a single extent, is used on the second coordinate. The component hom formula is
\[
 T\!\left(\coprod_a F(a)(x',x), B(b',b)\right)
 \cong\coprod_a T(F(a)(x',x),B(b',b)).
\]
Projection and cone-cell compatibility are those already proved in the product-index theorem. Thus separate fixed-interface preservation is available even when joint same-index preservation is not.

\paragraph{Relative inputs.} Start with the original tagged pure pseudofunctor $J:C\to K$, its derived directed parallel and compatible pure parallel, and the hypotheses of completed-pure-layer exactness. Apply local cocompletion and enriched module extension. Write $EJ:E(C)\to E(K)$. Interface objects $A,B$ are specified tagged images in $E(K)$, so their identity modules are pure. Define
\[
 P_A(X)=T_*(EJX,A),\qquad Q_B(X)=T_*(EJX,B)
\]
on the pure state bicategory $E(C)$. On composable pure modules, the parallel module compositor is invertible by whole-pure-layer exactness: each first-coordinate module is in the image of $EJ$, and each fixed second-coordinate identity is pure. Hence $P_A,Q_B$ are pseudofunctors, even though $T_*$ on arbitrary process pairs is only lax.

\paragraph{Fixed-interface state theorem.} Consider the native dialgebra bicategory whose objects are states $X$ with transitions $f:P_A X\to Q_B X$, whose arrows are $(h,\alpha:Q_Bh\,f\Rightarrow g\,P_Ah)$, and whose two-cells are compatible base two-cells. For a small pseudodiagram of these dialgebras, take its underlying pure-state lax cone object $A_F$ in $E(C)$. The functor $EJ$ preserves the explicit cone by its local coproduct/module construction. The fixed-interface calculation above then proves that $Q_B$ preserves it, with its actual projections. The native dialgebra cone-creation theorem therefore applies: the state cone carries a transition defined from the assembled component transitions, its projection squares are invertible, and it is universal for native cones and modifications. There is no comparison-adjoint hypothesis left to assume in this particular fixed-interface construction.

\paragraph{Proof dependency and dimensions.} This uses the previously proved native bicategorical dialgebra creation theorem, not a general claim that lax functors preserve bilimits. The pure state base $E(C)$ is a bicategory. Consequently the fixed-interface dialgebra is itself a bicategory; it is not automatically a Cat-valued hom of an ordinary bicategory of all state machines. Constructing the complete higher horizontal state structure, flattening and mixed parallel/feedback compatibility is additional work. The theorem supplies one necessary fixed-interface synchronization structure at its honest dimension.

\paragraph{Boundaries.} The result concerns the explicit normal-lax enriched cone presentations. It does not preserve every weighted bilimit, does not identify independent product indices with a shared diagonal, and does not extend parallel to all lax-action residuals. It does not turn retained-state feedback into instantaneous trace or solve general sharing. Its structural hypotheses are inherited from the original relative parallel construction plus its verified cocompletion/exactness conditions, rather than selected operational examples. No novelty claim is made about the underlying enriched module or dialgebra machinery.
\end{document}
